\documentclass[UTF8,fontset=fandol,11pt]{ctexart}
\usepackage[a4paper,margin=24mm]{geometry}
\usepackage{amsmath,amssymb,amsthm,graphicx,tabularx,booktabs,hyperref}
\newtheorem{definition}{定义}[section]
\newtheorem{theorem}[definition]{定理}
\newtheorem{example}[definition]{例题}
\newtheorem{exercise}[definition]{练习}
\hypersetup{unicode=true,colorlinks=true,linkcolor=blue}
\title{Offline draft recovery}
\date{}
\begin{document}
\maketitle
\par\noindent 离线合成演示：不代表真实模型对本专题的教学验收。\par
\section{Definition}\label{sec:S01}
% UNIT S01.U1
\begin{definition}
这里讨论实数的加法。
\end{definition}
% UNIT S01.U2
\begin{exercise}
\label{mathagent:exercise:9d1ffdd854167f1cbed066b927248fdfb8c3275b005db30f4a8cf5f1f6900586}
计算 $1+1$。
\end{exercise}
\section*{练习答案}
% UNIT S01.U2.answer
\par\noindent\textbf{\hyperref[mathagent:exercise:9d1ffdd854167f1cbed066b927248fdfb8c3275b005db30f4a8cf5f1f6900586]{练习 \ref*{mathagent:exercise:9d1ffdd854167f1cbed066b927248fdfb8c3275b005db30f4a8cf5f1f6900586} 的解答：}}$1+1=2$。\par
\end{document}
